% Part: lambda-calculus
% Chapter: church-rosser
% Section: definitions-and-properties

\documentclass[../../../include/open-logic-section]{subfiles}

\begin{document}

\olfileid{lam}{cr}{dap}

\olsection{నిర్వచనం మరియు లక్షణాలు}

ఈ అధ్యాయంలో చర్చ్--రోసర్ లక్షణం అనే భావాన్ని, దానికి
సంబంధించిన కొన్ని సాధారణ లక్షణాలను పరిచయం చేస్తాం.

\begin{defn}[చర్చ్--రోసర్ లక్షణం, CR] పదాలపై $\xredone$
  అనే సంబంధం
  \emph{చర్చ్--రోసర్ లక్షణాన్ని} తీరుస్తుందనడానికి
  అవసరమూ సరిపడేదీ ఇది: $M \xredone P$, $M \xredone Q$
  అయినప్పుడల్లా, $P \xredone N$, $Q \xredone N$
  అయ్యే ఏదో పదం~$N$ ఉండాలి.
\end{defn}

లాంబ్డా కలనశాస్త్రాన్ని ఒక గణన నమూనాగా చూడవచ్చు: అందులో
నియత రూపంలోని పదాలు “విలువలు”; పదాలపై తగ్గింపు సంబంధం
“గణన నియమాలు.” చర్చ్--రోసర్ లక్షణం చెప్పేది, ఒక గణనను
కొనసాగించడానికి ఒకటి కంటే ఎక్కువ మార్గాలున్నా, చివరి విలువ
లభించినట్లయితే అది ఒక్కటే అని.
\sourcecorrection{OLTELAMCRDAP-001}{మూలం ఈ మొదటి వివరణలో విలువ తప్పక ఉంటుందన్నట్లు చెబుతుంది. కానీ అదే విభాగం తరువాత “ఏదైనా ఉంటే” అని స్పష్టంగా షరతు పెడుతుంది; చర్చ్--రోసర్ లక్షణం ఒక్కటే ప్రతి పదం నియత రూపాన్ని చేరుతుందని చెప్పదు. తెలుగు వాక్యంలో విలువ లభిస్తే దాని అనన్యత మాత్రమే ప్రకటించాం.}

ప్రాథమిక బీజగణితం నుంచి ఉదాహరణగా $4 \times (1+2) + 3$ను
లెక్కించడానికి ఒకటి కంటే ఎక్కువ మార్గాలున్నాయి. ముందుగా
$1+2$ను~$3$గా తగ్గిస్తే $4 \times 3+3$ వస్తుంది;
లేదా పంపిణీ ధర్మంతో ముందుగా $4 \times (1+2)$ను
మారిస్తే $4 \times 1+4 \times 2+3$ వస్తుంది. ఈ రెండు
రూపాలనూ తరువాత $12+3$కు తగ్గించవచ్చు.

$\xredone$ను $\beta$-తగ్గింపుగా తీసుకుంటే, చర్చ్--రోసర్
లక్షణం వల్ల ఒక పదానికి నియత రూపం ఉంటే అది అనన్యమని
సులభంగా తెలుస్తుంది. $M$ పదం $P$, $Q$లకు తగ్గుతుందని,
అవి రెండూ నియత రూపాలనీ అనుకుందాం. చర్చ్--రోసర్ లక్షణం
ప్రకారం $P$, $Q$ రెండూ తగ్గే ఏదో $N$ ఉంది. కానీ
$P$, $Q$ నియత రూపాలే కాబట్టి $P$, $Q$ల నుంచి $N$కు
తగ్గింపు సున్నా దశలతోనే జరగాలి; అంటే $P$, $Q$, $N$
మూడూ ఒకే పదం. అందుకే ఒక పదం యొక్క
\emph{ఆ} నియత రూపం అని అంటాం.

లాంబ్డా కలనశాస్త్రాన్ని గణన నమూనాగా చూసినప్పుడు, ఒక
పదం నుంచి ప్రారంభమైన గణనకు దాని నియత రూపం “తుది
ఫలితం”గా భావించవచ్చు. పై పర్యవసానం ప్రకారం, తుది
ఫలితం ఏదైనా ఉంటే అది ఒక్కటే; అలాగే $4 \times (1+2)+3$
గణన ఫలితం ఒక్కటే, అది~$15$.

\begin{thm} \ollabel{thm:str}
  $\xredone$ సంబంధం చర్చ్--రోసర్ లక్షణాన్ని తీర్చి,
  $\xred$ అనేది $\xredone$ను కలిగి ఉండే కనిష్ఠ
  సంక్రమణ సంబంధమైతే, $\xred$ కూడా చర్చ్--రోసర్
  లక్షణాన్ని తీరుస్తుంది.
\end{thm}

\begin{proof}
  కింది రెండు తగ్గింపు క్రమాలున్నాయని అనుకుందాం:
  \begin{align*}
    M & \xredone P_1 \xredone \dots \xredone P_m \text{ మరియు}\\
    M & \xredone Q_1 \xredone \dots \xredone Q_n.
  \end{align*}
  ఎత్తు $m + 1$, వెడల్పు $n + 1$ గల పదాల $N$
  జాలకాన్ని నిర్మించి సిద్ధాంతాన్ని నిరూపిస్తాం. $i$వ
  అడ్డు వరుస, $j$వ నిలువు వరుసలోని పదాన్ని $N_{i,j}$తో
  సూచిస్తాం.

  $N_{i,j} \xredone N_{i+1,j}$,
  $N_{i,j} \xredone N_{i,j+1}$ అయ్యేలా $N$ను నిర్మిస్తాం.
  నిర్వచనం ఇది:
  \begin{align*}
    N_{0,0} &= M \\
    N_{i,0} &= P_i && \text{షరతు: } 1 \le i \le m \\
    N_{0,j} &= Q_j && \text{షరతు: } 1 \le j \le n \\
  \intertext{ఇతర సందర్భాల్లో:}
    N_{i,j} &= R
  \end{align*}
  ఇక్కడ $R$ అనేది $N_{i-1,j} \xredone R$,
  $N_{i,j-1} \xredone R$ షరతులను తీర్చే పదం.
  $\xredone$కు చర్చ్--రోసర్ లక్షణం ఉంది కాబట్టి
  అలాంటి పదం ఎల్లప్పుడూ ఉంటుంది.

  ఇప్పుడు $N_{m,0} \xredone \dots \xredone N_{m,n}$,
  $N_{0,n} \xredone \dots \xredone N_{m,n}$ ఉన్నాయి.
  $N_{m,0}$ అంటే $P_m$, $N_{0,n}$ అంటే $Q_n$.
  $\xred$ నిర్వచనం వల్ల సిద్ధాంతం వస్తుంది.
  \sourcecorrection{OLTELAMCRDAP-002}{మూల నిరూపణ చివర $N_{m,0}$ను $P$, $N_{0,n}$ను $Q$ అంటుంది; కానీ ఆ పేర్లు ముందుగా నిర్వచించలేదు. మూల క్రమాలూ జాలకపు సరిహద్దు నిర్వచనమూ వాటిని వరుసగా $P_m$, $Q_n$గా నిర్ధారిస్తున్నాయి. తెలుగు వాక్యంలో సూచికలతో ఉన్న ఆ రెండు తుది పదాలనే చూపాం; మూల క్రమాల గణిత సూత్రాలు మార్చలేదు.}
\end{proof}

\end{document}
